\chapter{\bilingual{İLGİLİ ÇALIŞMALAR}{LITERATURE REVIEW}}
\label{ch:literature}

% Citing references (MLA, via biblatex):
%   \autocite{key}                 -> (Author page)        parenthetical
%   \autocite[15]{key}             -> (Author 15)          with a page number
%   \autocite[pp.~15--18]{key}     -> (Author pp. 15-18)   with a page range
%   \textcite{key}                 -> Author (...)         name in running text
%   \autocites{a}{b}{c}            -> several works at once
% (Markdown/Pandoc syntax like [@key, p. 1] does NOT work in LaTeX.)
\bilingual{%
  Bu bölüm, tez konusuyla ilgili önceki çalışmaları özetler.
  Örneğin, kapsamlı bir LaTeX rehberi için bkz.\ \textcite{lamport1994};
  tipografik ilkeler için bkz.\ \autocite[ss.~23--28]{bringhurst2004}.}{%
  This chapter summarizes prior work related to the topic.
  For a comprehensive LaTeX reference see \textcite{lamport1994}; for
  typographic principles see \autocite[pp.~23--28]{bringhurst2004}.}

\section{\bilingual{Mevcut Şablonlar}{Existing Templates}}
\bilingual{%
  Birçok üniversite kendi tez şablonunu sağlar; ancak İstanbul Kültür
  Üniversitesi için resmi bir LaTeX şablonu bulunmamaktaydı
  \autocites{iku2004}{knuth1984}.}{%
  Many universities provide their own thesis templates; however, no official
  LaTeX template existed for Istanbul Kultur University
  \autocites{iku2004}{knuth1984}.}
